\bookchapter{}{Prelude: a drone in the wind}{The simulator, its three levels of reality, and the physics every lesson builds on.}
\label{ch:prelude}

\lead{Anemostatos — from the Greek \emph{ánemos}, wind, and \emph{statós}, standing: \enquote{the one that stands in the wind}. It is a simulator built for one purpose: to make feedback control visible. A quadcopter hovers above an endless grid, gravity pulls it down, random gusts push it around, and a controller fights back while every one of its terms is drawn on the screen. This chapter is the tour.}

\section*{The screen}

\screen[0 0 0 0]{meet}{The Anemostatos screen. Toolbar (top): run and pause, simulation speed, time, reset, the controller switch, a manual gust, the random seed, the level (L1/L2/L3) and a snapshot for comparing runs. 3D view (centre): the drone, the setpoint sphere, force arrows coloured by controller term, the live formula (bottom left) and the wind indicator (bottom right). Charts (bottom): tracking, error, controller terms, motor thrusts, and a switchable extra chart with the info card. Right: the lesson panel and every parameter, grouped.}{fig:prelude-screen}

\Cref{fig:prelude-screen} shows the whole application. Three things are worth knowing from the start.
\begin{itemize}
  \item \textbf{Everything is live.} Every slider acts immediately. Changing the controller type or the level restarts the run; everything else is changed in flight.
  \item \textbf{Everything is reproducible.} All randomness — wind, gusts, sensor noise — comes from one seeded generator. The same seed and the same parameters give exactly the same run, so two tunings can be compared on the same wind.\tip{When two runs differ, look at the seed first. With the same seed, every difference you see was caused by what you changed, not by luck.} The \ui{Snapshot} button (\key{S}) freezes the current run as a grey \emph{ghost} on the charts.
  \item \textbf{Everything is shareable.} The link button copies a URL that contains the complete parameter tree.
\end{itemize}

\section*{Three levels of reality}\index{simulator!levels of reality}

The simulator offers three models of the drone, of increasing fidelity. The lessons move from one to the next as the ideas require.

\begin{center}\small
\begin{tabularx}{\linewidth}{@{}l>{\raggedright\arraybackslash}X>{\raggedright\arraybackslash}X>{\raggedright\arraybackslash}X@{}}
\toprule
 & \sffamily\bfseries L1 · Altitude & \sffamily\bfseries L2 · Point mass & \sffamily\bfseries L3 · Quadrotor\\
\midrule
motion & vertical only & 3D translation & 3D translation + rotation\\
state & $y,\ \dot y$ & $\vec p,\ \vec v$ & $\vec p,\ \vec v,\ q,\ \vec\omega$, motors\\
input & total thrust $T$ & force vector $\vec F$ & four motor thrusts\\
controller & one PID & three PIDs & cascade of loops\\
lessons & 1–11, most of II.A–B & 12 & 12–14, most of Part II\\
\bottomrule
\end{tabularx}
\end{center}

At every level the drone is drawn as a full quadcopter with spinning propellers; only its freedom of motion changes. L2 is a deliberate cheat — real multicopters cannot push sideways — and \lessonref{les:tilt} exposes it.

\section*{The physics}\index{quadrotor!model}

The model is simple enough to reason about and rich enough to make control interesting. All units are SI; the world frame has $y$ pointing up and the floor at $y = 0$.

\subsection*{Motors}
Each motor $i$ receives a thrust command $f_i^{cmd} \in [0, f_{max}]$ with $f_{max} = \SI{6.1}{N}$, and its rotor responds with a first-order lag,
\begin{equation}\label{eq:prelude-motor}
  \tau_m\,\dot r_i = f_i^{cmd} - r_i,\qquad f_i = \eta_i\,r_i ,\qquad \tau_m = \SI{30}{ms}.
\end{equation}
$r_i$ is the rotor speed, expressed as the thrust a healthy propeller would make at that speed; the efficiency $\eta_i$ is 1 unless a lesson breaks a propeller (\lessonref{les:indi}). The four motors together lift at most \SI{24.4}{N}, a thrust-to-weight ratio of about 2.5 for the \SI{1}{kg} drone.\innumbers{A camera drone has a thrust-to-weight ratio near 2; a racing quadrotor 5 to 10. Below about 1.5 there is too little left over to fight a gust.}

\subsection*{Translation}
Newton's second law, with the thrust along the body's up-axis:
\begin{equation}\label{eq:prelude-trans}
  m\,\dot{\vec v} = R(q)\begin{pmatrix}0\\ \textstyle\sum_i f_i\\ 0\end{pmatrix} + m\vec g + \vec F_{drag} + \vec F_{ext}.
\end{equation}
$R(q)$ rotates body coordinates into world coordinates (it is the identity at L1). Drag acts on the velocity relative to the air, $\vec v_{rel} = \vec v - \vec w$, quadratically on each axis,
\begin{equation}
  F_{drag,j} = -c_j\,|v_{rel,j}|\,v_{rel,j},\qquad c_y = \SI{0.25}{}, \ c_x = c_z = \SI{0.12}{N.s^2/m^2},
\end{equation}
which is how the wind acts on the drone: a hovering drone in a wind $\vec w$ is pushed along the wind with a force that grows with its square.\innumbers{In a wind of \SI{4}{m/s} the hovering drone feels $0.12 \cdot 4^2 = \SI{1.9}{N}$ sideways: a fifth of its weight.}

\subsection*{Rotation (L3)}
The motors sit on an X-shaped frame (\cref{fig:prelude-motors}). Their thrust differences produce roll and pitch torques, and the rotors' aerodynamic drag produces a yaw torque, clockwise and counter-clockwise rotors in opposite directions:
\begin{equation}
  \tau_x = d\,(f_1 + f_4 - f_2 - f_3),\quad
  \tau_z = d\,(f_1 + f_2 - f_3 - f_4),\quad
  \tau_y = c_\tau\,(f_1 - f_2 + f_3 - f_4),
\end{equation}
with $d = L/\sqrt2$, $L = \SI{0.17}{m}$ and $c_\tau = \SI{0.016}{m}$. The body rotates according to Euler's equation and the attitude quaternion $q$ integrates the angular velocity:
\begin{equation}\label{eq:prelude-rot}
  J\,\dot{\vec\omega} = \vec\tau - \vec\omega \times (J\vec\omega) - c_\omega\vec\omega ,\qquad
  \dot q = \tfrac12\,q\otimes(0, \vec\omega).
\end{equation}
The inertia matrix is diagonal, $J = \diag(0.0082,\ 0.0149,\ 0.0082)\,\si{kg.m^2}$, the middle entry being the yaw axis. Quaternions — four numbers that describe a rotation without the singularities of angles — are explained in \lessonref{les:geometric} and the \hyperref[app:quat]{appendix}.
\marginnote{\resizebox{\marginparwidth}{!}{% Quadrotor X layout seen from above (+x forward, +z right), motor numbers and spin directions
\begin{tikzpicture}[scale=0.78, every node/.style={font=\sffamily\scriptsize}]
  \def\d{1.15}
  \draw[line width=2.2pt, ink!80, line cap=round] (-\d,-\d) -- (\d,\d) (-\d,\d) -- (\d,-\d);
  \fill[accent, rounded corners=2pt] (-0.35,-0.25) rectangle (0.35,0.25);
  \foreach \x/\y/\n/\dir in {\d/-\d/1/cw, \d/\d/2/ccw, -\d/\d/3/cw, -\d/-\d/4/ccw} {
    \draw[hair, fill=mist] (\y,\x) circle (0.55);
    \fill[ink] (\y,\x) circle (0.13);
    \node[text=ink] at ($(\y,\x)+(0,0.3)$) {M\n};
    \ifthenelse{\equal{\dir}{cw}}
      {\draw[->, sigP, line width=0.6pt] ($(\y,\x)+(150:0.42)$) arc (150:30:0.42);}
      {\draw[->, sigmeas, line width=0.6pt] ($(\y,\x)+(30:0.42)$) arc (30:150:0.42);}
  }
  \draw[->, line width=0.6pt] (0,1.95) -- (0,2.45) node[above] {$+x_b$ (front)};
  \draw[->, line width=0.6pt] (1.95,0) -- (2.45,0) node[right] {$+z_b$};
  \node[text=sigP] at (-2.3,-1.6) {CW};
  \node[text=sigmeas] at (-2.3,-1.9) {CCW};
\end{tikzpicture}
}\par\vspace{2pt}\captionsetup{type=figure}\caption{The motor layout seen from above.}\label{fig:prelude-motors}}[-134pt]

\subsection*{Ground}
The floor is solid. Touching it slowly means landing; touching it faster than \SI{3}{m/s}, or with a tilt of more than $80^\circ$, is a crash that stops the run. A run starts on the ground; the setpoint rises at \SI{1}{m/s} during take-off, and the integrators of the controllers are held at zero until the take-off is over.

\subsection*{Wind}
The air velocity is uniform in space and the sum of a constant mean wind, Ornstein–Uhlenbeck turbulence and discrete \enquote{one-minus-cosine} gusts. The details are in \lessonref{les:challenge}, where they matter most.

\subsection*{Sensors}
The controller never reads the true state. It reads a sensor model that can add Gaussian noise, a bias and a transport delay to each signal, and can hold the position between slow fixes (like a barometer or a GPS receiver). Besides position, velocity, attitude and body rates, the sensors provide an accelerometer — which measures the \emph{specific force} $R^\T(\dot{\vec v} - \vec g)$, so it reads $+g$ upward while hovering and zero in free fall — and the rotor speeds of the motors.\didyouknow{A phone finds \enquote{down} this way: at rest its accelerometer reads $g$, pointing straight up.} Part~II uses both.

\section*{Time}

The physics advances in fixed steps of \SI{1}{ms} with the semi-implicit (symplectic) Euler method: velocities are updated from the forces, then positions from the \emph{new} velocities;\pitfall{The order matters. With the plain Euler method — positions from the \emph{old} velocities — an undamped oscillation gains a little energy at every step and slowly explodes.} the same for the angular velocity and the quaternion, which is renormalised every step. The controllers run on their own clocks — \SI{250}{Hz} for the altitude controller, up to \SI{1}{kHz} for the fastest loop of the quadrotor — and hold their outputs between samples, exactly like the software of a real flight controller. Rendering is decoupled from all of this: the screen shows interpolated states at whatever frame rate the browser manages.

\begin{inapp}[The source code]
The simulator is written in TypeScript. The physics and the controllers are plain classes with no dependence on the user interface, so they can run headlessly — which is how this book's data were produced.
\begin{itemize}[leftmargin=1.2em]
  \item \src{src/sim/} — the drone, the dynamics (\texttt{dynamics.ts}), wind (\texttt{wind.ts}), sensors (\texttt{sensors.ts}), keep-out zones, and all parameters with their defaults (\texttt{params.ts}).
  \item \src{src/control/} — every controller of the book, one file each, from \texttt{pid.ts} to \texttt{policy.ts}.
  \item \src{src/estimation/} — the Kalman filter and the disturbance observers of Part~II.
  \item \src{src/engine/simulation.ts} — the loop that ties it together; \texttt{telemetry.ts} records every signal at \SI{200}{Hz}.
  \item \src{src/lessons/} — the lessons of this book, as code: set-ups, scripted events, goals.
\end{itemize}
\end{inapp}

\clearpage
\section*{Notation and conventions}
\label{sec:notation}

The same symbol means the same thing in every chapter. Where the literature uses one letter for two things, this book keeps them apart, and the table says how.

\begin{fullwidth}
\begin{center}\small
\renewcommand{\arraystretch}{1.15}
\begin{tabularx}{\linewidth}{@{}l>{\raggedright\arraybackslash}X@{\hspace{1.2em}}l>{\raggedright\arraybackslash}X@{}}
\toprule
\multicolumn{4}{@{}l}{\sffamily\bfseries Signals of a loop}\\
$r$ & setpoint (reference) & $y$ & controlled variable; at L1 the altitude\\
$e = r - y$ & error, positive when too low & $x = y - r$ & deviation from the setpoint, $x = -e$\\
$u$ & controller output (command) & $d$ & disturbance, positive upwards\\
$T$ & total thrust & $f_i$ & thrust of motor $i$\\
\midrule
\multicolumn{4}{@{}l}{\sffamily\bfseries The controller and the plant}\\
$K_p,\ K_i,\ K_d$ & PID gains & $I$ & the integral term, $I = K_i\!\int\! e\,\dd t$\\
$m,\ \hat m$ & true and assumed mass & $g$ & \SI{9.81}{m/s^2}\\
$\tau_m$ & motor time constant & $c_v$ & vertical drag coefficient\\
$\Delta t$ & controller sample period & $t_k = k\,\Delta t$ & sampling instants; $e_k = e(t_k)$\\
$\omega_n,\ \zeta$ & natural frequency, damping ratio & $\omega_d$ & damped frequency\\
$T_{osc}$ & period of an oscillation & $\tau$ & a time constant or a delay\\
\midrule
\multicolumn{4}{@{}l}{\sffamily\bfseries Vectors, matrices, complex numbers}\\
$\vx,\ \symbf{u}$ & state and input vectors: bold & $A,\ B,\ C,\ K$ & matrices: italic capitals\\
$\Id$ & the identity matrix & $A^\T$ & transpose\\
$\vec p,\ \vec v,\ \vec\omega$ & vectors in 3D space: arrows & $J$ & inertia (a number or a matrix)\\
$q,\ R(q)$ & attitude quaternion, rotation matrix (body to world) & $\|\cdot\|,\ |\cdot|$ & Euclidean norm, absolute value\\
$s = \sigma + j\omega$ & complex frequency; $j = \sqrt{-1}$ & $\operatorname{Re},\ \operatorname{Im}$ & real and imaginary part\\
\bottomrule
\end{tabularx}
\end{center}
\end{fullwidth}

\begin{description}[leftmargin=0pt, itemsep=1pt, topsep=2pt, font=\sffamily\small\bfseries]
  \item[Marks.] A dot is a derivative with respect to time, $\dot y = \dd y/\dd t$. A hat is an estimate or a belief ($\hat m$, $\hat{\vx}$). A star is an equilibrium ($e^*$). The subscript $ss$ is a steady-state value, $sp$ a setpoint passed from one loop to the next, $0$ an initial value.
  \item[Frames and signs.] The world frame has $y$ pointing up and $x$, $z$ horizontal. Forces, positions and velocities along $y$ are positive upwards. The body frame is fixed to the drone, with its $y$-axis along the thrust.
  \item[Units.] SI throughout. Angles are in radians and frequencies in \si{rad/s} in every formula; degrees and hertz appear only where they are written out.\pitfall{Hertz and radians per second differ by $2\pi \approx 6.3$. A \SI{5}{Hz} oscillation is a \SI{31}{rad/s} oscillation. Mixing the two is a classic source of wrong numbers.}
  \item[Statements.] A \emph{Definition} fixes a term. A \emph{Theorem} is a general statement with its hypotheses, followed by a proof, whose end is marked \rule{0.45em}{0.45em}, or by the source where the proof is found. A \emph{Result} is a formula derived for our drone. A \emph{Key idea} is the same thing in words. Theorems are numbered within each lesson, and the sources are listed in the \hyperref[app:refs]{references}.
\end{description}
